\documentclass[twoside]{article}

% AISTATS 2027 had not released an author kit when this working draft was
% assembled. We use the latest official AISTATS style without modification.
\usepackage{aistats2026}

\usepackage[round,authoryear]{natbib}
\bibliographystyle{apalike}
\renewcommand{\bibname}{References}
\renewcommand{\bibsection}{\subsubsection*{\bibname}}

\usepackage{amsmath,amssymb,amsthm,mathtools,bm}
\usepackage{booktabs,array,multirow,graphicx,tabularx}
\usepackage{placeins}
\usepackage{microtype}
\usepackage{xcolor}
\usepackage[hidelinks]{hyperref}
\usepackage{url}

\input{results_placeholders}

\newcommand{\E}{\mathbb{E}}
\newcommand{\Pp}{\mathbb{P}}
\newcommand{\Rr}{\mathbb{R}}
\newcommand{\Var}{\operatorname{Var}}
\newcommand{\Cov}{\operatorname{Cov}}
\newcommand{\tr}{\operatorname{tr}}
\newcommand{\diag}{\operatorname{diag}}
\newcommand{\argmin}{\operatorname*{arg\,min}}
\newcommand{\cF}{\mathcal{F}}
\newcommand{\cE}{\mathcal{E}}
\newcommand{\cR}{\mathcal{R}}
\newcommand{\cX}{\mathcal{X}}
\newcommand{\1}{\mathbf{1}}
\newcommand{\pos}[1]{[#1]_{+}}
\newcommand{\dif}{\mathop{}\!\mathrm{d}}
\newcommand{\Normal}{\mathcal{N}}
\newcommand{\Err}{\operatorname{Err}}

\theoremstyle{plain}
\newtheorem{theorem}{Theorem}[section]
\newtheorem{proposition}[theorem]{Proposition}
\newtheorem{corollary}[theorem]{Corollary}
\newtheorem{lemma}[theorem]{Lemma}
\theoremstyle{definition}
\newtheorem{assumption}[theorem]{Assumption}
\theoremstyle{remark}
\newtheorem{remark}[theorem]{Remark}

\begin{document}

\runningtitle{When Is One Flow Enough?}
\runningauthor{Anonymous Authors}

\twocolumn[
\aistatstitle{When Is One Flow Enough?\\
Identifiability and Model Selection for Opposing-Flow Dynamics}
\aistatsauthor{Anonymous Authors}
\aistatsaddress{Anonymous Institution}
]

\begin{abstract}
Many observed transitions are net effects of two opposing processes. We ask when a signed one-flow model is sufficient and when the processes should be represented by two nonnegative conditional flows. For the drift $m_0(x,y)=U_0(x)(1-y)-R_0(x)y$, separate identification holds if and only if the conditional state variance is positive. The exact population loss from suppressing one flow is $G(x)=v(x)\kappa(x)$, where $v(x)$ measures state dispersion and $\kappa(x)$ measures coactivity. In a local Gaussian benchmark, these factors combine into $\Gamma_n=nG_n/\sigma_\varepsilon^2$; the two-flow fit has lower expected prediction error precisely when $\Gamma_n>1$. An exactly solvable two-state design displays the crossover and separates identification from model-selection value. We instantiate both classes with parameter-matched neural dynamics and evaluate them on complete held-out storm events. On complete held-out storm events, the two-flow model reduces one-step error by \MainStepGain\ and 24-hour forecast error by \MainForecastGain, with the largest gains where turnover information and coactivity are jointly high. The conclusion is conditional: separating opposing flows is useful when its representation benefit exceeds the statistical cost of estimating the additional component.
\end{abstract}

\section{Introduction}

County-level outage records reveal the fraction of customers without service, not the underlying counts of newly interrupted and newly restored customers. Utility records can retrospectively decompose resilience curves into overlapping outage and restoration processes \citep{carrington2021}; in contrast, a prominent statistical outage model treats the simultaneous movements through a single net process because separate modeling is considered impractical \citep{zhu2021}. These positions pose a statistical question that is more precise than whether two processes can be named: when does their separation improve prediction from aggregate state observations?

We study the bounded-state transition
\begin{equation}
Y_{t+1}=Y_t+U_0(X_t)(1-Y_t)-R_0(X_t)Y_t+\varepsilon_t,
\label{eq:intro_transition}
\end{equation}
where $U_0$ is an interruption component, $R_0$ is a restoration component, and $X_t$ denotes observed drivers. The competing one-flow representation uses one signed function: positive values act on the served pool $1-Y_t$, negative values act on the interrupted pool $Y_t$. It can switch direction across driver values, but it cannot keep both directions positive at the same driver value. The two-flow representation is therefore more expressive, but it also estimates an additional function. The relevant question is not whether two flows are possible; it is whether their representation gain repays their estimation cost.

The analysis has one organizing quantity. In a fixed-design Gaussian benchmark, let $G_n$ denote the approximation loss of the best one-flow span and let $\sigma_\varepsilon^2$ be the one-step noise variance. The ratio $\Gamma_n=nG_n/\sigma_\varepsilon^2$ factors into information for total turnover and the cost of suppressing one active direction. This factorization unifies three issues that are often discussed separately: whether two flows are needed, whether they can be identified, and whether the additional component is worth estimating.

Our contributions follow this chain. First, we derive the exact population loss of the state-scaled one-flow class, including nonnegative-boundary and degenerate-state cases. Second, an orthogonal net-drift/turnover parameterization gives an exact information decomposition and an explicit representation--estimation criterion; a supplied exactly solvable case reduces the criterion to one scalar crossover. Third, we report a single real-data study in which the two neural classes differ only in flow separation, together with the strongest same-information baseline and a simple operational baseline. The empirical result is deliberately reported as mixed: the structural contrast is favorable on average at 24 hours but does not pass the original confirmatory gate, and neither neural flow class dominates the practical baselines at long horizons.

At fixed $x$, Equation~\eqref{eq:intro_transition} is a varying-coefficient regression \citep{hastie1993}. We do not claim the general construction as novel. The contribution is the exact geometry generated by the basis $(1-y,-y)$ and its implication for choosing between one signed flow and two nonnegative flows. Neural differential models provide flexible parameterizations of unknown dynamics \citep{chen2018}; here the neural networks instantiate the two model classes rather than constitute a separate architectural contribution.

\input{sections/02_related_work}
\input{sections/03_setup}
\input{sections/04_theory}
\input{sections/05_experiments}
\input{sections/06_discussion}

\bibliography{references}

\clearpage
\input{sections/07_checklist}

\clearpage
\appendix
\onecolumn
\aistatstitle{When Is One Flow Enough?\\Supplementary Material}
\input{appendices/A_proofs}
\input{appendices/B_finite_sample}
\input{appendices/C_transfer_rollout}
\input{appendices/D_experimental_protocol}
\input{appendices/E_supplementary_results}
\input{appendices/F_numerical_audit}

\end{document}
